% ============================================================
% CARGA ELÉCTRICA Y ELECTRIZACIÓN
% ============================================================


% ============================================================
% INTRODUCCIÓN AL FENÓMENO
% ============================================================

\section{Carga eléctrica y electrización}

Antes de hablar de campos eléctricos, potenciales o circuitos, debemos
responder una pregunta más fundamental:

\begin{quote}
¿Qué propiedad de la materia hace posibles los fenómenos eléctricos?
\end{quote}

Todos hemos observado alguna manifestación de la electricidad.

En un día seco, al quitarnos una prenda podemos escuchar pequeños
chasquidos. Después de caminar sobre una alfombra podemos sentir una
descarga al tocar una superficie metálica. Un peine frotado contra el cabello
puede atraer pequeños trozos de papel. En una escala muchísimo mayor,
durante una tormenta observamos descargas eléctricas en forma de rayos.

Aunque estos fenómenos parecen diferentes, todos están relacionados con una
propiedad fundamental de la materia: la \textbf{carga eléctrica}.

\begin{notapedagogica}
Una forma útil de comenzar la clase es realizar o plantear una experiencia
muy sencilla.

Frote un peine plástico, una regla o un globo contra el cabello y acérquelo
a pequeños trozos de papel.

Antes de explicar el fenómeno, pregunte:

\begin{enumerate}
    \item ¿Por qué el plástico puede atraer el papel después de frotarlo?
    \item ¿Apareció algo nuevo en el plástico?
    \item ¿Se creó carga eléctrica al frotarlo?
    \item Si el papel es eléctricamente neutro, ¿por qué puede ser atraído?
\end{enumerate}

No es necesario responder inmediatamente todas las preguntas. Conviene
utilizarlas como hilo conductor de la clase.
\end{notapedagogica}


% ============================================================
% ELECTROSTÁTICA
% ============================================================

\subsection{Electrostática}

El estudio de la electricidad suele dividirse inicialmente según el estado
de las cargas que analizamos.

\begin{definicion}[Electrostática]
La \textbf{electrostática} es la parte de la Física que estudia las
propiedades e interacciones de distribuciones de carga eléctrica que pueden
considerarse en reposo.
\end{definicion}

Esto no significa que las cargas nunca puedan haberse movido.

Por ejemplo, durante un proceso de electrización pueden transferirse
electrones de un cuerpo a otro. Después de esa transferencia podemos estudiar
la distribución resultante mediante la electrostática.

Más adelante estudiaremos cargas en movimiento de forma sostenida, lo cual
nos conducirá al concepto de corriente eléctrica.


% ============================================================
% LA MATERIA Y SU ESTRUCTURA ELÉCTRICA
% ============================================================

\subsection{La materia posee estructura eléctrica}

Para comprender la carga eléctrica debemos recordar que la materia está
formada por átomos.

En una descripción elemental, un átomo contiene:

\begin{itemize}
    \item protones;
    \item neutrones;
    \item electrones.
\end{itemize}

Los protones y neutrones se encuentran en el núcleo, mientras que los
electrones forman la estructura electrónica que rodea al núcleo.

Estas partículas presentan diferentes propiedades eléctricas.

\begin{table}[H]
    \centering

    \caption{Carga eléctrica de las principales partículas del átomo.}

    \begin{tabular}{lcc}
        \toprule
        Partícula & Símbolo & Carga eléctrica \\
        \midrule
        Protón    & \(p\) & \(+e\) \\
        Electrón  & \(e^{-}\) & \(-e\) \\
        Neutrón   & \(n\) & \(0\) \\
        \bottomrule
    \end{tabular}

\end{table}

La magnitud \(e\) recibe el nombre de \textbf{carga elemental}.

\begin{equation}
    e
    =
    \SI{1.602176634e-19}{\coulomb}.
    \label{eq:carga_elemental}
\end{equation}

La carga del protón es

\begin{equation}
    q_p=+e,
\end{equation}

mientras que la carga del electrón es

\begin{equation}
    q_e=-e.
\end{equation}

Las dos poseen la misma magnitud y signos opuestos.

\begin{idea}
En los procesos electrostáticos ordinarios de sólidos no suelen desplazarse
los protones de un objeto a otro.

Los protones permanecen ligados a los núcleos.

Las partículas que normalmente se transfieren o redistribuyen son los
\textbf{electrones}.
\end{idea}


% ============================================================
% DEFINICIÓN DE CARGA ELÉCTRICA
% ============================================================

\subsection{¿Qué es la carga eléctrica?}

\begin{definicion}[Carga eléctrica]
La \textbf{carga eléctrica} es una propiedad física fundamental de la materia
que determina su participación en las interacciones electromagnéticas.
\end{definicion}

Representaremos habitualmente la carga mediante la letra

\[
q
\]

o, cuando nos refiramos a la carga total de un sistema,

\[
Q.
\]

La carga eléctrica es una \textbf{magnitud escalar}.

Esto significa que distintas cantidades de carga se combinan
algebraicamente teniendo en cuenta sus signos.

Si un sistema contiene varias cargas,

\begin{equation}
    Q=q_1+q_2+\cdots+q_n.
    \label{eq:carga_total}
\end{equation}


% ============================================================
% UNIDAD
% ============================================================

\subsection{Unidad de carga eléctrica}

En el Sistema Internacional, la unidad de carga eléctrica es el
\textbf{coulomb}, cuyo símbolo es

\[
\si{\coulomb}.
\]

El coulomb es una cantidad de carga muy grande comparada con la carga de un
solo electrón.

Como

\[
e=\SI{1.602176634e-19}{\coulomb},
\]

un coulomb corresponde aproximadamente a

\begin{equation}
    \frac{\SI{1}{\coulomb}}{e}
    \approx
    6.24\times10^{18}
\end{equation}

cargas elementales.

\begin{interpretacion}
Esto ayuda a entender por qué una carga macroscópica aparentemente pequeña,
por ejemplo algunos microcoulombs, representa en realidad un enorme número
de electrones transferidos.
\end{interpretacion}


% ============================================================
% DOS TIPOS DE CARGA
% ============================================================

\subsection{Carga positiva y carga negativa}

Experimentalmente encontramos dos clases de carga eléctrica, que por
convención llamamos:

\begin{itemize}
    \item \textbf{carga positiva};
    \item \textbf{carga negativa}.
\end{itemize}

El protón posee carga positiva y el electrón carga negativa.

La elección de los nombres positivo y negativo es histórica. Podrían haberse
intercambiado los nombres sin modificar la Física, siempre que toda la
convención se utilizara consistentemente.


% ============================================================
% INTERACCIÓN CUALITATIVA
% ============================================================

\subsection{Interacción entre cargas}

Una de las primeras propiedades observadas experimentalmente es que los
cuerpos cargados pueden atraerse o repelerse.

La regla cualitativa es:

\begin{equation}
    \boxed{
    \text{cargas del mismo signo se repelen}
    }
\end{equation}

y

\begin{equation}
    \boxed{
    \text{cargas de signos opuestos se atraen}
    }.
\end{equation}

Por tanto:

\begin{itemize}
    \item positiva con positiva: repulsión;
    \item negativa con negativa: repulsión;
    \item positiva con negativa: atracción.
\end{itemize}

\begin{notapedagogica}
En esta clase conviene quedarse todavía en la descripción cualitativa.

La pregunta

\begin{quote}
¿Cuánto vale la fuerza de atracción o repulsión?
\end{quote}

se responderá en el siguiente capítulo mediante la ley de Coulomb.

Aquí lo fundamental es establecer primero qué es la carga y cuáles son sus
propiedades.
\end{notapedagogica}


% ============================================================
% CUERPO NEUTRO
% ============================================================

\subsection{¿Qué significa que un cuerpo sea eléctricamente neutro?}

Un error frecuente consiste en pensar que un cuerpo neutro no posee cargas
eléctricas.

Esto no es correcto.

Un cuerpo ordinario contiene una enorme cantidad de protones y electrones.

\begin{definicion}[Cuerpo eléctricamente neutro]
Un cuerpo es \textbf{eléctricamente neutro} cuando su carga eléctrica neta
es cero.
\end{definicion}

Si un cuerpo contiene igual cantidad total de carga positiva y negativa,

\[
Q_{+}+Q_{-}=0,
\]

entonces

\[
Q=0.
\]

\begin{idea}
\textbf{Neutro no significa ``sin cargas''.}

Significa que las cargas positivas y negativas se compensan exactamente en
la carga neta.
\end{idea}


% ============================================================
% EXCESO Y DEFECTO DE ELECTRONES
% ============================================================

\subsection{¿Cómo adquiere carga un cuerpo?}

En los fenómenos electrostáticos ordinarios, un cuerpo se carga
principalmente porque cambia su número de electrones.

Si un cuerpo inicialmente neutro \textbf{pierde electrones}, queda con un
exceso relativo de carga positiva.

Por tanto,

\begin{equation}
    \boxed{
    \text{pierde electrones}
    \Longrightarrow
    \text{carga positiva}
    }.
\end{equation}

Si un cuerpo inicialmente neutro \textbf{gana electrones}, posee un exceso de
carga negativa.

Por tanto,

\begin{equation}
    \boxed{
    \text{gana electrones}
    \Longrightarrow
    \text{carga negativa}
    }.
\end{equation}

\begin{notapedagogica}
Conviene detenerse aquí y preguntar:

\begin{quote}
Si un cuerpo queda cargado positivamente,
¿ganó cargas positivas?
\end{quote}

En los procesos de electrización ordinarios, la respuesta es
\textbf{no necesariamente y, en sólidos comunes, normalmente no}.

Generalmente el cuerpo quedó positivo porque
\textbf{perdió electrones}.
\end{notapedagogica}

\begin{errorfrecuente}
No decir que un cuerpo metálico se carga positivamente porque
``ganó protones''.

Los protones permanecen ligados en los núcleos.

En los procesos electrostáticos habituales entre cuerpos sólidos, son los
electrones los que se transfieren.
\end{errorfrecuente}


% ============================================================
% EXPRESIÓN MEDIANTE NÚMERO DE ELECTRONES
% ============================================================

\subsection{Carga neta y número de electrones transferidos}

Si un cuerpo inicialmente neutro pierde \(N\) electrones, su carga es

\begin{equation}
    Q=+Ne.
    \label{eq:carga_perdida_electrones}
\end{equation}

Si gana \(N\) electrones,

\begin{equation}
    Q=-Ne.
    \label{eq:carga_ganada_electrones}
\end{equation}

Por tanto, si conocemos la magnitud de la carga podemos determinar cuántos
electrones han sido transferidos:

\begin{equation}
    \boxed{
    N=\frac{\abs{Q}}{e}
    }.
    \label{eq:numero_electrones_carga}
\end{equation}


% ============================================================
% EJEMPLO
% ============================================================

\subsection{Ejemplo: cuerpo que pierde electrones}

\begin{ejemplo}
Un cuerpo inicialmente neutro pierde

\[
N=5.0\times10^{12}
\]

electrones.

¿Cuál es su carga final?
\end{ejemplo}

Como el cuerpo \textbf{pierde electrones}, sabemos antes de hacer cualquier
cálculo que deberá quedar con carga positiva.

Utilizamos

\[
Q=+Ne.
\]

Entonces,

\begin{align}
    Q
    &=
    \paren{5.0\times10^{12}}
    \paren{\SI{1.602176634e-19}{\coulomb}}
    \\
    &\approx
    \SI{8.0e-7}{\coulomb}.
\end{align}

Por tanto,

\begin{equation}
    \boxed{
    Q\approx+\SI{8.0e-7}{\coulomb}
    }.
\end{equation}

También puede escribirse como

\begin{equation}
    \boxed{
    Q\approx+\SI{0.80}{\micro\coulomb}
    }.
\end{equation}

\begin{interpretacion}
El signo positivo no apareció por una regla matemática arbitraria.

Lo anticipamos físicamente desde el comienzo: el cuerpo perdió electrones,
por lo que quedó con déficit de carga negativa.
\end{interpretacion}


% ============================================================
% CONSERVACIÓN DE LA CARGA
% ============================================================

\section{Conservación de la carga eléctrica}

Una de las propiedades fundamentales de la carga eléctrica es su
conservación.

\begin{principio}[Principio de conservación de la carga eléctrica]
La carga eléctrica total de un sistema aislado permanece constante.
\end{principio}

Matemáticamente,

\begin{equation}
    \boxed{
    Q_{\text{total,inicial}}
    =
    Q_{\text{total,final}}
    }.
    \label{eq:conservacion_carga}
\end{equation}

Esto significa que en los procesos ordinarios de electrización la carga no
aparece de la nada ni desaparece.

Lo que ocurre es una \textbf{transferencia} o una
\textbf{redistribución} de carga.


% ============================================================
% EJEMPLO CON DOS CUERPOS
% ============================================================

\subsection{Interpretación de la conservación}

Supongamos que tenemos dos cuerpos inicialmente neutros.

Por tanto,

\[
Q_{\text{inicial}}=0.
\]

Después de cierto proceso, el cuerpo A pierde electrones y queda con carga

\[
+Q.
\]

Los electrones transferidos llegan al cuerpo B, que queda con carga

\[
-Q.
\]

La carga total final es

\begin{equation}
    Q_{\text{final}}
    =
    (+Q)+(-Q)
    =
    0.
\end{equation}

Así,

\[
Q_{\text{inicial}}
=
Q_{\text{final}}.
\]

\begin{idea}
Cuando dos cuerpos inicialmente neutros se electrizan entre sí y el sistema
puede considerarse aislado, uno no puede adquirir una carga sin que aparezca
una carga compensadora en otra parte del sistema.
\end{idea}


% ============================================================
% FROTAMIENTO Y CONSERVACIÓN
% ============================================================

\begin{notapedagogica}
Este es un momento apropiado para regresar al experimento inicial del peine,
la regla o el globo.

Puede preguntarse:

\begin{quote}
Cuando frotamos dos cuerpos, ¿el frotamiento crea electrones?
\end{quote}

La respuesta es no.

Los electrones ya existían.

El proceso favorece su transferencia de un material al otro.

Por eso la electrización por frotamiento debe explicarse como
\textbf{transferencia de carga}, no como creación de carga.
\end{notapedagogica}


% ============================================================
% CUANTIZACIÓN
% ============================================================

\section{Cuantización de la carga eléctrica}

La carga eléctrica posee otra propiedad fundamental: está cuantizada.

\begin{principio}[Cuantización de la carga]
La carga eléctrica neta observable en los procesos ordinarios aparece en
múltiplos enteros de la carga elemental \(e\).
\end{principio}

Por tanto,

\begin{equation}
    \boxed{
    Q=ne
    },
    \qquad
    n=0,\pm1,\pm2,\pm3,\ldots
    \label{eq:cuantizacion_carga}
\end{equation}

Esto significa que un cuerpo puede tener, por ejemplo,

\[
Q=e,
\]

\[
Q=-3e,
\]

o

\[
Q=1500e,
\]

pero en el contexto electrostático ordinario no encontraremos un objeto
aislado con

\[
Q=1.37e.
\]

\begin{interpretacion}
La palabra \textit{cuantizada} significa que la carga no cambia de manera
arbitrariamente continua a escala microscópica.

La transferencia ocurre en unidades elementales asociadas, en los procesos
ordinarios, al número entero de electrones transferidos.
\end{interpretacion}


% ============================================================
% EJEMPLO DE CUANTIZACIÓN
% ============================================================

\subsection{Ejemplo: identificar una carga posible}

Supongamos que una partícula tiene una carga

\[
Q=-3.204353268\times10^{-19}\,\si{\coulomb}.
\]

Dividimos su magnitud entre la carga elemental:

\begin{align}
    N
    &=
    \frac{\abs{Q}}{e}
    \\
    &=
    \frac{3.204353268\times10^{-19}}
    {1.602176634\times10^{-19}}
    \\
    &=2.
\end{align}

Por tanto,

\[
Q=-2e.
\]

La carga corresponde a dos cargas elementales negativas.

\begin{nota}
En la física de partículas existen constituyentes fundamentales, como los
quarks, a los que se asignan cargas fraccionarias de \(e\).

Sin embargo, los quarks no se observan aislados en las situaciones
electrostáticas ordinarias. Para el estudio que desarrollamos aquí, la
relación \(Q=ne\) describe las cargas netas que encontraremos.
\end{nota}


% ============================================================
% CONDUCTORES Y AISLANTES
% ============================================================

\section{Conductores y aislantes}

Hasta ahora hemos preguntado cuánto vale la carga de un cuerpo.

Ahora debemos hacer una pregunta diferente:

\begin{quote}
¿Qué tan fácilmente puede moverse la carga dentro del material?
\end{quote}

La respuesta depende de la estructura microscópica del material.


% ============================================================
% CONDUCTORES
% ============================================================

\subsection{Conductores}

\begin{definicion}[Conductor eléctrico]
Un \textbf{conductor} es un material en el cual existen portadores de carga
que pueden desplazarse con relativa facilidad a través del material.
\end{definicion}

Los metales son ejemplos importantes de conductores.

En un metal, algunos electrones se encuentran débilmente ligados a átomos
individuales y pueden desplazarse a través de la estructura metálica.

Por esta razón, cuando depositamos carga en un conductor, esta puede
redistribuirse.

Entre los conductores encontramos, por ejemplo:

\begin{itemize}
    \item cobre;
    \item aluminio;
    \item plata;
    \item otros metales;
    \item el cuerpo humano, en muchas situaciones;
    \item soluciones que contienen iones móviles.
\end{itemize}

\begin{idea}
Que un material sea conductor no significa que tenga más electrones que un
aislante.

La diferencia fundamental está en la
\textbf{movilidad de los portadores de carga}.
\end{idea}


% ============================================================
% AISLANTES
% ============================================================

\subsection{Aislantes}

\begin{definicion}[Aislante eléctrico]
Un \textbf{aislante} es un material en el cual las cargas no pueden
desplazarse libremente a través de grandes distancias dentro del material en
las condiciones consideradas.
\end{definicion}

En los aislantes, los electrones se encuentran mucho más ligados a átomos o
moléculas particulares.

Son ejemplos comunes:

\begin{itemize}
    \item plástico;
    \item vidrio;
    \item caucho;
    \item cerámica;
    \item aire seco en condiciones ordinarias.
\end{itemize}

Si frotamos una barra aislante y depositamos carga en una región, esa carga
puede permanecer relativamente localizada en ella.

Esto contrasta con un conductor, donde la carga tiene mayor libertad para
redistribuirse.


% ============================================================
% COMPARACIÓN
% ============================================================

\subsection{Comparación entre conductor y aislante}

\begin{table}[H]
    \centering

    \caption{Comparación cualitativa entre conductores y aislantes.}

    \begin{tabularx}{\textwidth}{lYY}
        \toprule
        & \textbf{Conductor} & \textbf{Aislante} \\
        \midrule

        Movilidad de carga
        &
        Relativamente alta
        &
        Muy reducida
        \\

        Redistribución
        &
        La carga puede redistribuirse fácilmente
        &
        La carga suele permanecer más localizada
        \\

        Ejemplos
        &
        Metales
        &
        Plástico, vidrio, caucho
        \\

        \bottomrule
    \end{tabularx}

\end{table}

\begin{advertencia}
La distinción entre conductor y aislante no es absolutamente rígida.

La conductividad depende del material y también de condiciones como
temperatura, humedad, impurezas y magnitud del campo eléctrico.

Por ejemplo, el aire normalmente actúa como aislante, pero puede ionizarse
durante una descarga eléctrica intensa y permitir el paso de carga.
\end{advertencia}


% ============================================================
% POLARIZACIÓN
% ============================================================

\section{Polarización eléctrica}

Existe una situación que parece contradictoria:

\begin{quote}
¿Cómo puede un objeto cargado atraer a un cuerpo que es eléctricamente
neutro?
\end{quote}

La respuesta está en la redistribución de las cargas.


% ============================================================
% POLARIZACIÓN EN CONDUCTOR
% ============================================================

\subsection{Polarización de un conductor}

Consideremos un conductor neutro.

Inicialmente su carga neta es cero.

Ahora acercamos, sin tocarlo, una barra con carga negativa.

Los electrones móviles del conductor experimentan repulsión y se desplazan
hacia la región más alejada de la barra.

Como consecuencia:

\begin{itemize}
    \item en la región próxima a la barra aparece un déficit relativo de
    electrones;

    \item en la región más alejada aparece un exceso de electrones.
\end{itemize}

El conductor continúa teniendo carga neta cero, pero sus cargas se han
separado espacialmente.

\begin{definicion}[Polarización]
La \textbf{polarización eléctrica} es la separación o redistribución de
cargas positivas y negativas dentro de un cuerpo, sin que necesariamente
cambie su carga neta.
\end{definicion}

\begin{idea}
Un objeto puede ser:

\[
\text{neutro}
\]

y al mismo tiempo

\[
\text{polarizado}.
\]

No existe contradicción.

Ser neutro se refiere a la \textbf{carga neta}.

Estar polarizado se refiere a la \textbf{distribución espacial} de las
cargas.
\end{idea}


% ============================================================
% ATRACCIÓN DE UN CUERPO NEUTRO
% ============================================================

\subsection{¿Por qué una carga puede atraer un objeto neutro?}

Supongamos que un cuerpo negativo se acerca a un objeto neutro.

La carga del objeto neutro puede redistribuirse o polarizarse de modo que las
cargas de signo opuesto queden, en promedio, más cerca del cuerpo cargado.

Como la interacción eléctrica depende de la distancia, la atracción de las
cargas más cercanas puede superar la repulsión de las cargas más alejadas.

El resultado puede ser una fuerza neta de atracción.

\begin{notapedagogica}
Aquí puede retomarse el experimento del peine y los trozos de papel.

La conclusión importante es:

\begin{quote}
Que un cuerpo cargado atraiga a otro cuerpo no demuestra por sí solo que
ambos posean cargas netas de signos opuestos.
\end{quote}

Un cuerpo eléctricamente neutro también puede ser atraído debido a
polarización.
\end{notapedagogica}

\begin{errorfrecuente}
La regla

\begin{quote}
``si dos objetos se atraen, entonces tienen cargas opuestas''
\end{quote}

no siempre es válida.

La repulsión sí permite concluir que ambos objetos tienen cargas netas del
mismo signo.

La atracción puede ocurrir entre cargas opuestas o entre un cuerpo cargado
y otro neutro polarizado.
\end{errorfrecuente}


% ============================================================
% ELECTRIZACIÓN
% ============================================================

\section{Electrización}

Ya podemos responder una pregunta central:

\begin{quote}
¿Qué significa electrizar un cuerpo?
\end{quote}

\begin{definicion}[Electrización]
La \textbf{electrización} es el proceso mediante el cual un cuerpo adquiere
una carga eléctrica neta o modifica su estado eléctrico mediante transferencia
o redistribución de cargas.
\end{definicion}

En los fenómenos ordinarios estudiaremos tres métodos principales:

\begin{enumerate}
    \item electrización por frotamiento;
    \item electrización por contacto;
    \item electrización por inducción.
\end{enumerate}


% ============================================================
% FROTAMIENTO
% ============================================================

\subsection{Electrización por frotamiento}

Consideremos dos cuerpos inicialmente neutros hechos de materiales diferentes.

Al ponerlos en contacto y frotarlos, determinados electrones pueden
transferirse de un material al otro.

El cuerpo que pierde electrones queda cargado positivamente.

El cuerpo que gana electrones queda cargado negativamente.

Si el sistema puede considerarse aislado,

\[
Q_{\text{total,inicial}}=0
\]

y, después del proceso,

\[
Q_1+Q_2=0.
\]

Por tanto,

\begin{equation}
    Q_1=-Q_2.
\end{equation}

\begin{idea}
El frotamiento no crea carga eléctrica.

Facilita la transferencia de electrones entre los materiales.
\end{idea}

\begin{notapedagogica}
La explicación en clase puede formularse de manera muy directa:

\begin{quote}
Antes de frotarlos, los dos cuerpos son neutros.

Durante el contacto algunos electrones pasan de uno al otro.

El que pierde electrones queda positivo.

El que gana electrones queda negativo.

Si consideramos ambos cuerpos juntos, la carga total sigue siendo la misma.
\end{quote}

De esta manera se conecta inmediatamente la electrización con el principio
de conservación de la carga.
\end{notapedagogica}


% ============================================================
% CONTACTO
% ============================================================

\subsection{Electrización por contacto}

Supongamos ahora que tenemos un conductor cargado y otro conductor
inicialmente neutro.

Cuando los ponemos en contacto, los portadores de carga pueden desplazarse
entre ambos cuerpos.

La carga se redistribuye hasta alcanzar una nueva situación de equilibrio.

\begin{definicion}[Electrización por contacto]
La \textbf{electrización por contacto} ocurre cuando un cuerpo adquiere
carga eléctrica mediante contacto directo con otro cuerpo cargado.
\end{definicion}

Después de separarlos, el cuerpo que inicialmente era neutro puede quedar
cargado con el mismo signo que el cuerpo que produjo la electrización.

\begin{advertencia}
No debe suponerse automáticamente que la carga se divide exactamente por la
mitad.

Eso ocurre únicamente en situaciones especiales, por ejemplo cuando se
ponen en contacto dos conductores idénticos y aislados bajo condiciones
simétricas.

En general, la distribución final depende de la geometría y de las
características de los conductores.
\end{advertencia}


% ============================================================
% INDUCCIÓN
% ============================================================

\subsection{Electrización por inducción}

Existe un método aún más interesante.

Es posible cargar un conductor sin que el cuerpo utilizado para producir la
electrización llegue a tocarlo.

Este procedimiento recibe el nombre de \textbf{inducción}.

\begin{definicion}[Electrización por inducción]
La \textbf{electrización por inducción} es un proceso mediante el cual un
conductor adquiere carga neta debido a la influencia de un cuerpo cargado,
sin contacto directo con este.
\end{definicion}


% ============================================================
% CONEXIÓN A TIERRA
% ============================================================

\subsection{Conexión a tierra}

Para comprender la inducción debemos introducir la idea de
\textbf{conexión a tierra}.

La Tierra puede actuar, en muchas situaciones electrostáticas, como un enorme
reservorio capaz de recibir o suministrar electrones sin experimentar un
cambio apreciable en su estado eléctrico global.

Representaremos la conexión a tierra mediante un camino conductor entre el
objeto y la Tierra.


% ============================================================
% PROCEDIMIENTO DE INDUCCIÓN
% ============================================================

\subsection{Proceso de carga por inducción}

Supongamos que queremos cargar un conductor inicialmente neutro utilizando
una barra con carga negativa.

El procedimiento puede analizarse paso a paso.

\subsubsection*{Paso 1: acercar la barra}

Acercamos la barra negativa al conductor sin tocarlo.

Los electrones móviles del conductor son repelidos hacia el lado opuesto.

El conductor se polariza.

Todavía,

\[
Q_{\text{conductor}}=0.
\]


\subsubsection*{Paso 2: conectar el conductor a tierra}

Mientras la barra permanece cerca, conectamos el conductor a tierra.

Debido a la repulsión producida por la barra negativa, algunos electrones
pueden abandonar el conductor y pasar hacia la Tierra.

Ahora el conductor ha perdido electrones.


\subsubsection*{Paso 3: retirar primero la conexión a tierra}

Sin retirar todavía la barra cargada, desconectamos el conductor de la
Tierra.

Los electrones que abandonaron el conductor ya no pueden regresar por ese
camino.


\subsubsection*{Paso 4: retirar la barra inductora}

Finalmente alejamos la barra negativa.

La carga que permanece en el conductor se redistribuye.

Como el conductor perdió electrones, queda con carga neta positiva.

\begin{equation}
    \boxed{
    \text{inductor negativo}
    \Longrightarrow
    \text{conductor inducido positivo}
    }
\end{equation}

Si utilizamos un inductor positivo, el proceso se invierte y el conductor
puede quedar cargado negativamente.

\begin{idea}
En una electrización por inducción bien realizada, el cuerpo inducido queda
con signo opuesto al del cuerpo inductor.
\end{idea}


% ============================================================
% DIFERENCIA ENTRE POLARIZACIÓN E INDUCCIÓN
% ============================================================

\subsection{Polarización no es lo mismo que electrización por inducción}

Esta distinción es fundamental.

Si simplemente acercamos una barra cargada a un conductor neutro,

\[
\text{se redistribuyen las cargas},
\]

pero

\[
Q_{\text{neto}}=0.
\]

El conductor está \textbf{polarizado}, pero sigue siendo neutro.

Para que quede finalmente cargado por inducción debe existir un mecanismo
que permita transferencia de carga, por ejemplo una conexión a tierra.

\begin{errorfrecuente}
No llamar automáticamente ``carga por inducción'' a cualquier separación de
cargas causada por acercar un objeto cargado.

Acercar el objeto produce inicialmente
\textbf{polarización}.

La electrización por inducción requiere que el cuerpo termine con una carga
neta diferente de cero.
\end{errorfrecuente}


% ============================================================
% COMPARACIÓN DE LOS MÉTODOS
% ============================================================

\subsection{Comparación de los métodos de electrización}

\begin{table}[H]
    \centering

    \caption{Comparación cualitativa de los principales métodos de
    electrización.}

    \begin{tabularx}{\textwidth}{lYYY}
        \toprule
        Método
        &
        ¿Hay contacto?
        &
        ¿Hay transferencia de electrones?
        &
        Resultado típico
        \\
        \midrule

        Frotamiento
        &
        Sí
        &
        Sí
        &
        Los cuerpos quedan con cargas de signos opuestos
        \\

        Contacto
        &
        Sí
        &
        Sí
        &
        El cuerpo neutro adquiere normalmente el signo del cuerpo cargado
        \\

        Inducción
        &
        No con el inductor
        &
        Sí, mediante tierra u otro camino conductor
        &
        El cuerpo inducido queda con signo opuesto al inductor
        \\

        \bottomrule
    \end{tabularx}

\end{table}


% ============================================================
% UNA DISTINCIÓN MUY IMPORTANTE
% ============================================================

\subsection{Transferencia de carga y redistribución de carga}

A lo largo de esta clase han aparecido dos procesos diferentes.

\subsubsection{Transferencia}

En una transferencia, electrones pasan de un cuerpo a otro.

Esto ocurre, por ejemplo, durante:

\begin{itemize}
    \item frotamiento;
    \item contacto;
    \item inducción cuando existe conexión a tierra.
\end{itemize}


\subsubsection{Redistribución}

En una redistribución, las cargas que ya se encuentran dentro de un cuerpo
cambian de posición sin que necesariamente cambie la carga neta del cuerpo.

Esto ocurre durante la polarización.

\begin{idea}
Esta distinción será muy importante más adelante:

\[
\boxed{
\text{redistribuir carga}
\neq
\text{crear carga}
}
\]

y

\[
\boxed{
\text{polarizar un cuerpo}
\neq
\text{darle necesariamente carga neta}
}.
\]
\end{idea}


% ============================================================
% ERRORES FRECUENTES
% ============================================================

\section{Errores y confusiones frecuentes}

\begin{errorfrecuente}
\textbf{Pensar que un cuerpo neutro no contiene cargas.}

Un cuerpo neutro posee cargas positivas y negativas; simplemente su suma
algebraica es cero.
\end{errorfrecuente}


\begin{errorfrecuente}
\textbf{Decir que la fricción crea carga eléctrica.}

La electrización por frotamiento implica transferencia de electrones.

La carga total se conserva.
\end{errorfrecuente}


\begin{errorfrecuente}
\textbf{Pensar que un cuerpo positivo ganó protones.}

En los procesos electrostáticos ordinarios de sólidos, un cuerpo queda
positivo normalmente porque perdió electrones.
\end{errorfrecuente}


\begin{errorfrecuente}
\textbf{Confundir carga positiva con ``más carga''.}

Los signos positivo y negativo indican dos tipos de carga, no una comparación
entre cantidades grandes y pequeñas.
\end{errorfrecuente}


\begin{errorfrecuente}
\textbf{Pensar que los electrones tienen mayor carga que los protones.}

La magnitud de las cargas es la misma:

\[
\abs{q_e}
=
\abs{q_p}
=
e.
\]

Sus signos son opuestos.
\end{errorfrecuente}


\begin{errorfrecuente}
\textbf{Confundir conductor con cuerpo cargado.}

Un conductor puede ser eléctricamente neutro.

Ser conductor describe la movilidad de los portadores de carga, no la carga
neta del cuerpo.
\end{errorfrecuente}


\begin{errorfrecuente}
\textbf{Confundir aislante con material sin cargas eléctricas.}

Los aislantes también están formados por partículas cargadas.

La diferencia es que sus cargas tienen mucha menor libertad de movimiento
macroscópico.
\end{errorfrecuente}


\begin{errorfrecuente}
\textbf{Concluir que dos cuerpos que se atraen necesariamente poseen cargas
opuestas.}

Un cuerpo cargado puede atraer un cuerpo neutro debido a polarización.
\end{errorfrecuente}


\begin{errorfrecuente}
\textbf{Confundir polarización con electrización por inducción.}

La polarización redistribuye las cargas.

La electrización por inducción deja finalmente al cuerpo con una carga neta.
\end{errorfrecuente}


% ============================================================
% SECUENCIA PARA DICTAR LA CLASE
% ============================================================

\section{Secuencia conceptual para explicar el tema}

\begin{notapedagogica}
Una secuencia especialmente clara para dictar esta clase es:

\begin{enumerate}

    \item Comenzar con una experiencia:
    peine, regla o globo frotado que atrae pequeños trozos de papel.

    \item Preguntar qué cambió en el cuerpo después de frotarlo.

    \item Introducir la idea de electrostática.

    \item Recordar brevemente la estructura del átomo:
    protones, neutrones y electrones.

    \item Introducir la carga eléctrica como propiedad fundamental de la
    materia.

    \item Establecer los dos signos de carga:
    positiva y negativa.

    \item Explicar cualitativamente:

    \[
    \text{iguales se repelen},
    \qquad
    \text{opuestas se atraen}.
    \]

    \item Explicar qué significa realmente que un cuerpo sea neutro.

    \item Establecer la idea central:

    \begin{quote}
    perder electrones produce carga positiva;
    ganar electrones produce carga negativa.
    \end{quote}

    \item Introducir el coulomb y la carga elemental.

    \item Explicar el principio de conservación de la carga.

    \item Volver al frotamiento y aclarar que la carga no fue creada:
    simplemente se transfirió.

    \item Introducir la cuantización:

    \[
    Q=ne.
    \]

    \item Explicar conductores y aislantes desde la movilidad de los
    portadores de carga.

    \item Mostrar que un objeto neutro puede polarizarse.

    \item Utilizar la polarización para explicar por qué un cuerpo cargado
    puede atraer objetos neutros.

    \item Explicar electrización por frotamiento.

    \item Explicar electrización por contacto.

    \item Explicar electrización por inducción paso a paso, incluyendo la
    conexión a tierra.

    \item Comparar los tres métodos.

    \item Finalizar conectando el tema con la siguiente pregunta:

    \begin{quote}
    Ya sabemos que las cargas pueden atraerse o repelerse.

    ¿Cómo podemos calcular la magnitud de esa fuerza?
    \end{quote}

    Esta pregunta conduce directamente a la ley de Coulomb.

\end{enumerate}
\end{notapedagogica}


% ============================================================
% PREGUNTAS DURANTE LA CLASE
% ============================================================

\section{Preguntas para comprobar comprensión}

Antes de pasar a ejercicios formales pueden realizarse preguntas orales como:

\begin{enumerate}

    \item ¿Un cuerpo eléctricamente neutro carece de cargas eléctricas?

    \item ¿Cuál es la diferencia entre un cuerpo neutro y un cuerpo cargado?

    \item ¿Qué ocurre si un cuerpo inicialmente neutro pierde electrones?

    \item ¿Qué ocurre si gana electrones?

    \item ¿Por qué no decimos normalmente que un objeto se carga
    positivamente porque ganó protones?

    \item ¿Qué tienen en común la magnitud de la carga del protón y la del
    electrón?

    \item ¿Qué significa afirmar que la carga se conserva?

    \item Cuando frotamos dos cuerpos neutros, ¿se está creando carga?

    \item ¿Qué significa que la carga eléctrica está cuantizada?

    \item ¿Puede existir, dentro del modelo estudiado, una carga neta de

    \[
    2.5e?
    \]

    \item ¿Cuál es la diferencia fundamental entre un conductor y un
    aislante?

    \item ¿Puede un conductor ser eléctricamente neutro?

    \item ¿Puede un aislante contener electrones?

    \item ¿Puede un cuerpo cargado atraer a otro cuerpo eléctricamente
    neutro?

    \item Si dos cuerpos se repelen, ¿qué podemos afirmar acerca de sus
    cargas?

    \item ¿Cuál es la diferencia entre polarización y electrización?

    \item ¿En cuál método de electrización no es necesario tocar el cuerpo
    con el objeto cargado?

    \item ¿Qué función cumple la conexión a tierra durante una electrización
    por inducción?

\end{enumerate}


% ============================================================
% RESUMEN
% ============================================================

\begin{resumen}

\begin{itemize}

    \item La carga eléctrica es una propiedad fundamental de la materia
    relacionada con las interacciones electromagnéticas.

    \item Existen dos signos de carga:

    \[
    \text{positiva}
    \qquad\text{y}\qquad
    \text{negativa}.
    \]

    \item Cargas del mismo signo se repelen y cargas de signo opuesto se
    atraen.

    \item El protón tiene carga

    \[
    q_p=+e,
    \]

    el electrón

    \[
    q_e=-e,
    \]

    y el neutrón tiene carga neta cero.

    \item La carga elemental es

    \[
    e=
    \SI{1.602176634e-19}{\coulomb}.
    \]

    \item La unidad SI de carga es el coulomb.

    \item Un cuerpo neutro no carece de cargas; posee carga neta igual a
    cero.

    \item En los procesos electrostáticos ordinarios:

    \[
    \text{perder electrones}
    \Longrightarrow
    Q>0,
    \]

    mientras que

    \[
    \text{ganar electrones}
    \Longrightarrow
    Q<0.
    \]

    \item La carga eléctrica se conserva:

    \[
    Q_{\text{inicial}}
    =
    Q_{\text{final}}
    \]

    para un sistema aislado.

    \item La carga está cuantizada:

    \[
    Q=ne.
    \]

    \item Un conductor permite que los portadores de carga se desplacen con
    relativa facilidad.

    \item En un aislante las cargas tienen mucha menor libertad para
    desplazarse a través del material.

    \item La polarización consiste en una redistribución de cargas y puede
    ocurrir sin cambiar la carga neta del cuerpo.

    \item Un objeto cargado puede atraer un objeto neutro debido a
    polarización.

    \item Los principales mecanismos de electrización estudiados son:

    \begin{enumerate}
        \item frotamiento;
        \item contacto;
        \item inducción.
    \end{enumerate}

    \item La electrización por frotamiento implica transferencia de
    electrones entre materiales.

    \item La electrización por contacto implica transferencia directa de
    carga al poner cuerpos en contacto.

    \item La electrización por inducción permite cargar un conductor sin
    contacto directo con el cuerpo inductor y normalmente requiere un camino
    para transferir carga, como una conexión a tierra.

    \item La carga no se crea durante estos procesos: se transfiere o se
    redistribuye.

\end{itemize}

\end{resumen}